% Chapter 4: Justification
\section{Justificativa}

A relevância deste projeto reside em três pilares principais: educacional, tecnológico e social.

Do ponto de vista educacional, este trabalho contribui com uma nova ferramenta didática que aborda um desafio persistente no ensino de química. A Tabela Periódica é a base para a compreensão de inúmeros conceitos químicos, e facilitar seu aprendizado pode ter um impacto positivo significativo no desempenho dos alunos na disciplina. A abordagem de aprendizagem baseada em jogos está alinhada com teorias construtivistas modernas, que defendem que o conhecimento é construído ativamente pelo aprendiz \cite{journal_example}.

Tecnologicamente, o projeto explora a aplicação de tecnologias de desenvolvimento de jogos para a criação de soluções educacionais, uma área em plena expansão. O desenvolvimento de "Alchemon" servirá como um estudo de caso prático sobre o design e a implementação de jogos sérios.

Socialmente, a democratização do acesso a ferramentas educacionais de qualidade por meio de dispositivos móveis pode ajudar a reduzir disparidades no aprendizado. Um jogo como "Alchemon" pode ser utilizado tanto em sala de aula, como um complemento às atividades do professor, quanto em casa, como uma forma de estudo autônomo.
